\documentclass[11pt,letterpaper]{article}

\usepackage[margin=1in]{geometry}
\usepackage{amsmath,amssymb,amsthm,mathtools}
\usepackage{authblk}
\usepackage{hyperref}
\usepackage{microtype}
\usepackage{enumitem}

\numberwithin{equation}{section}

\theoremstyle{plain}
\newtheorem{theorem}{Theorem}[section]
\newtheorem{lemma}[theorem]{Lemma}
\newtheorem{proposition}[theorem]{Proposition}
\newtheorem{corollary}[theorem]{Corollary}
\theoremstyle{definition}
\newtheorem{definition}[theorem]{Definition}
\newtheorem{example}[theorem]{Example}
\theoremstyle{remark}
\newtheorem{remark}[theorem]{Remark}

\newcommand{\A}{\mathcal{A}}
\newcommand{\D}{\mathcal{D}}
\newcommand{\NN}{\mathbb{N}}
\newcommand{\ZZ}{\mathbb{Z}}
\newcommand{\Xinf}{X_\infty}
\newcommand{\supp}{\mathrm{supp}}
\DeclareMathOperator{\Val}{Val}
\DeclareMathOperator{\Fold}{Fold}
\DeclareMathOperator{\lcp}{lcp}

\renewcommand\Authfont{\normalsize}
\renewcommand\Affilfont{\small}
\setlength{\affilsep}{0.4em}

\title{Canonical Zeckendorf Normalization and the Minimal Berstel Adder}
\author[1]{Haobo Ma}
\author[2]{Wenlin Zhang\thanks{Corresponding author.}}
\affil[1]{AELF PTE LTD., \#14-02, Marina Bay Financial Centre Tower 1, 8 Marina Blvd, Singapore 018981, Singapore\\\texttt{auric@aelf.io}}
\affil[2]{National University of Singapore, Singapore\\\texttt{e1327962@u.nus.edu}}
\date{}

\begin{document}
\maketitle

\begin{abstract}
We study the canonical normalization map behind addition in Zeckendorf
numeration.  Digitwise superposition of two admissible Zeckendorf words yields
a finitely supported word over $\{0,1,2\}$, and addition is recovered by
replacing this overlay with the unique admissible word of the same Fibonacci
value.  We realize this normalization as the canonical section of a Fibonacci
congruence and describe it by a terminating local rewrite system whose
irreducible configurations are exactly the admissible Zeckendorf words.  We
then separate this one-sided low-to-high canonical normalizer from the
classical most-significant-digit-first adders of Frougny and Berstel.  In the
one-sided model, the canonical map has coordinatewise unbounded anticipation,
is not subsequential, and has exact prefix-destruction index $\Delta(n)=n$.
By contrast, the classical Berstel transducer belongs to a different
subsequential setting, and we prove that its state complexity is exactly $10$.
Finally, we show that every fixed-radius synchronous local realization of
canonical normalization requires a number of rounds growing linearly with the
input scale.
\end{abstract}

\medskip
\noindent\textbf{Keywords.}
Zeckendorf numeration, canonical normalization, subsequential transducer,
Berstel adder, online delay, Fibonacci addition.

\medskip
\noindent\textbf{Mathematics Subject Classification.}
68Q45, 68Q70, 11B39.

\section{Background and Main Results}

Zeckendorf addition can be viewed as a normalization problem.  If
$c,d\in\Xinf$ are admissible Zeckendorf words, then their digitwise overlay
$c+d$ is a finitely supported word over $\{0,1,2\}$, and the arithmetic task
is to replace this noncanonical word by the unique admissible word with the
same Fibonacci value.  Thus the central object of the paper is the canonical
normalization map
\[
\Fold_\infty:\{0,1,2\}^{(\NN)}\to \Xinf.
\]
Our purpose is to describe this map in structural terms and to compare two
different finite-state viewpoints on Fibonacci addition.

Finite-state arithmetic in nonstandard numeration systems is classical,
especially in Pisot bases
\cite{Frougny1991FibonacciFiniteAutomata,Frougny1992NumbersAutomata,Frougny1999OnlineAddition}.  Recent work
continues to treat Fibonacci-based numeration from an automata-theoretic
viewpoint; see, for instance,
\cite{AhlbachUsatineFrougnyPippenger2013,LabbeLepsova2023FibonacciTwosComplement}.
For broader finite-state normalization phenomena in linear and Ostrowski-type
numeration systems, see also
\cite{Frougny1992NumbersAutomata,HieronymiTerry2018OstrowskiAddition}.  The
present note remains strictly Fibonacci-specific and does not attempt PLRS,
$k$-bonacci, or multidimensional generalizations.
The point of departure here is that the canonical map $\Fold_\infty$ on
one-sided low-to-high streams should not be conflated with the classical
most-significant-digit-first adders.  We therefore begin with the
normalization mechanism itself.  First, we realize $\Fold_\infty$ as the
canonical section of the Fibonacci congruence on the free commutative monoid
of scale counts and give a terminating local rewrite system whose irreducible
configurations are exactly the admissible Zeckendorf words.  This yields the
factorization
\[
\text{superpose digits} \longrightarrow \Fold_\infty
\]
for Zeckendorf addition.

The first complexity statement is negative, but its novelty is task-specific
rather than historical.  Earlier work of Frougny already showed that
finite-state behavior in Fibonacci numeration depends delicately on the
direction of reading and on the chosen conversion problem
\cite{Frougny1991FibonacciFiniteAutomata,Frougny1992NumbersAutomata,Frougny1999OnlineAddition}.
In particular, finite-state normalization results in the literature often pass
through compositions of one-sided subsequential transformations rather than a
single canonical low-to-high map \cite{Frougny1992NumbersAutomata}.  Our
negative result concerns exactly that single canonical map.
What we add here is a sharp formulation for the canonical one-sided
normalization map itself.  In the low-to-high model natural for
$\Fold_\infty$, every fixed output coordinate has unbounded anticipation.
Equivalently, the map admits no finite one-pass online delay.  More strongly,
once finitely supported inputs are encoded as low-to-high finite words, the
resulting normalization function is not subsequential.  The obstruction is not
a minor technicality: for every fixed lookahead bound, there are two inputs
agreeing on an arbitrarily long prefix whose selected normalized digits
nevertheless differ, and one may even choose a prefix pair for which the first
normalized digit flips.  In fact, for every prescribed output length $n$ there
is a prefix-related input pair whose normalized outputs disagree from the very
first symbol while the shorter output has length exactly $n$.  This sharply
separates canonical one-sided
normalization from the classical online-addition constructions of Frougny,
where the input is read most significant digit first and the output is not the
greedy expansion \cite[Prop.~13]{Frougny1999OnlineAddition}.  The witnesses
may already be chosen inside the genuine addition image $\{c+d:c,d\in\Xinf\}$.

We then revisit the classical $10$-state Berstel adder in its proper setting.
The existence of this transducer, and a $10$-state presentation of it, are
classical
\cite{Frougny1999OnlineAddition,LabbeLepsova2023FibonacciTwosComplement}.  Our
new point is an exact minimality statement.  Using deterministic
subsequential Nerode equivalence, we show that this classical transducer is
already minimal: every deterministic subsequential realization of the same
recoding requires at least $10$ reachable states.

This sequential distinction should be compared with the literature on parallel
addition in nonstandard numeration systems
\cite{FrougnyPelantovaSvobodova2013MinimalDigitSets}.  Our final result shows
that the canonical map $\Fold_\infty$ remains intrinsically nonlocal in the
synchronous setting as well: every fixed-radius local normalization scheme
requires a number of rounds growing linearly with the input scale.

The paper is organized as follows.  Section~\ref{sec:preliminaries} introduces
the quotient description of Zeckendorf normalization.  Section~\ref{sec:overlay}
gives the rewrite system and identifies addition with overlay plus
normalization.  Section~\ref{sec:streaming} proves that the one-sided map
$\Fold_\infty$ has coordinatewise unbounded anticipation, no bounded online
delay, and no subsequential realization in the low-to-high finite-word model.
Sections~\ref{sec:kernel} and \ref{sec:minimality} turn to the classical
Berstel adder and prove that its deterministic subsequential state complexity
is exactly $10$.  Section~\ref{sec:complexity} contains the local parallel
lower bound.

\section{Canonical Normalization in Zeckendorf Numeration}\label{sec:preliminaries}

\subsection{Admissible Words and the Value Map}

\paragraph{Fibonacci weights.}
Throughout the paper we use the convention
\[
F_1=F_2=1,\qquad F_{n+2}=F_{n+1}+F_n\quad (n\ge 1).
\]
The weight at position $k$ is $F_{k+1}$.

\paragraph{Admissible representatives.}
Let
\[
\Xinf
:=
\left\{
c=(c_1,c_2,\ldots)\in\{0,1\}^{\NN} :
\supp(c)\text{ is finite and } c_kc_{k+1}=0\text{ for all }k
\right\}.
\]
For $c\in \Xinf$ we define
\[
\Val(c):=\sum_{k\ge 1} c_k F_{k+1}\in \NN.
\]
By Zeckendorf's theorem, every $N\in\NN$ admits a unique representative
$Z(N)\in \Xinf$ such that $\Val(Z(N))=N$ \cite{Zeckendorf1972}.

\paragraph{Free scale counts.}
Let
\[
\D
:=
\left\{
a=(a_1,a_2,\ldots) : a_k\in\NN,\ \supp(a)\text{ finite}
\right\},
\]
with pointwise addition.  We write $e_k$ for the $k$th standard basis vector.

\begin{definition}[Fibonacci congruence]\label{def:fib-congruence}
Let $\equiv$ be the least commutative-monoid congruence on $\D$ generated by
\[
e_2\equiv 2e_1
\qquad\text{and}\qquad
e_{k+2}\equiv e_{k+1}+e_k \quad (k\ge 1).
\]
\end{definition}

\begin{definition}[Value map on scale counts]\label{def:val-on-D}
For $a\in\D$ we define
\[
\Val(a):=\sum_{k\ge 1} a_k F_{k+1}.
\]
\end{definition}

\begin{proposition}[Completeness of Fibonacci congruence]\label{prop:val-invariant}
If $a\equiv b$ in $\D$, then $\Val(a)=\Val(b)$.  Consequently the value map
descends to a well-defined monoid homomorphism
\[
\overline{\Val}:\D/\!\equiv \to \NN.
\]
\end{proposition}

\begin{proof}
It is enough to check the defining relations.  For the recurrence relation,
\[
\Val(e_{k+2})=F_{k+3}=F_{k+2}+F_{k+1}=\Val(e_{k+1}+e_k).
\]
For the base relation,
\[
\Val(e_2)=F_3=2=2F_2=\Val(2e_1).
\]
The claim follows because $\equiv$ is generated from these relations by
additivity.
\end{proof}

\begin{theorem}[The quotient monoid is $\NN$]\label{thm:monoid-quotient-is-N}
The homomorphism $\overline{\Val}:\D/\!\equiv \to \NN$ is an isomorphism of
commutative monoids.  Moreover every congruence class contains a unique
representative in $\Xinf$.
\end{theorem}

\begin{proof}
Surjectivity is immediate from Zeckendorf's theorem: for every $N\in\NN$,
$Z(N)\in \Xinf$ satisfies $\Val(Z(N))=N$.  For injectivity, write
$\varepsilon=[e_1]\in\D/\!\equiv$.  By induction on the recurrence relation,
\[
[e_k]=F_{k+1}\varepsilon \qquad (k\ge 1).
\]
Hence every class $[a]\in\D/\!\equiv$ satisfies
\[
[a]
=
\sum_{k\ge 1} a_k[e_k]
=
\left(\sum_{k\ge 1} a_kF_{k+1}\right)\varepsilon
=
\Val(a)\varepsilon.
\]
So the congruence class is uniquely determined by its value.  The uniqueness
of the representative in $\Xinf$ is then exactly the uniqueness of the
Zeckendorf expansion.
\end{proof}

\subsection{Canonical Representatives and the Fold Section}

\begin{corollary}[The fold map is a computable section]\label{cor:fold-as-section}
The embedding $\Xinf\hookrightarrow \D$ is a computable section of the quotient
map $\D\twoheadrightarrow \D/\!\equiv \cong \NN$.  Equivalently, the map
\[
\Fold_\infty:\D\to \Xinf,\qquad
\Fold_\infty(a):=Z(\Val(a)),
\]
is a well-defined normalization map selecting the unique admissible
representative in each congruence class.
\end{corollary}

\begin{proof}
The statement is a restatement of Theorem~\ref{thm:monoid-quotient-is-N}.
\end{proof}

\section{Overlay and Normalization}\label{sec:overlay}

\subsection{Local Value-Preserving Rewrites}

Although $\Fold_\infty$ is defined quotient-theoretically, its action can be
described directly by local digit rewrites on $\D$.  For
$a=(a_1,a_2,\ldots)\in\D$, consider the following directed moves:
\begin{align*}
\text{(R1)}\qquad
a&\to a-e_k-e_{k+1}+e_{k+2}
&&\text{if } a_k\ge 1,\ a_{k+1}\ge 1,\\
\text{(R2)}\qquad
a&\to a-2e_1+e_2
&&\text{if } a_1\ge 2,\\
\text{(R3)}\qquad
a&\to a-2e_2+e_1+e_3
&&\text{if } a_2\ge 2,\\
\text{(R4)}\qquad
a&\to a-2e_k+e_{k-2}+e_{k+1}
&&\text{if } k\ge 3,\ a_k\ge 2.
\end{align*}
Each rule preserves value, since it is a directed form of one of the
identities
\[
F_{k+1}+F_{k+2}=F_{k+3},\qquad
2F_2=F_3,\qquad
2F_3=F_2+F_4,\qquad
2F_{k+1}=F_{k-1}+F_{k+2}\ (k\ge 3).
\]

\begin{proposition}[Termination and unique normal form]\label{prop:rewrite-termination}
Starting from any $a\in\D$, every rewrite sequence generated by
\textup{(R1)}--\textup{(R4)} terminates.  The irreducible configurations are
exactly the admissible words $\Xinf$, and the terminal irreducible word is
$\Fold_\infty(a)$.
\end{proposition}

\begin{proof}
Define
\[
A(a):=\sum_{k\ge 1} a_k,\qquad
B(a):=\sum_{k\ge 1} k\,a_k,\qquad
C(a):=a_2,
\]
and compare triples $\mu(a):=(A(a),B(a),C(a))$ lexicographically.
Rule \textup{(R1)} and rule \textup{(R2)} decrease $A$ by $1$.
Rule \textup{(R4)} preserves $A$ and decreases $B$ by $1$.
Rule \textup{(R3)} preserves both $A$ and $B$, but decreases $C$ by $2$.
Hence every rewrite step strictly decreases $\mu(a)$, so infinite rewrite
sequences are impossible.

If $a$ is irreducible, then no adjacent pair of positive digits can occur, or
else \textup{(R1)} would apply; likewise no digit can be $\ge 2$, or else one of
\textup{(R2)}--\textup{(R4)} would apply.  Thus every irreducible configuration is a
finitely supported $\{0,1\}$-word with no adjacent $1$s, i.e. an element of
$\Xinf$.  Conversely, no rule applies to an element of $\Xinf$, so the
irreducibles are exactly $\Xinf$.

Every rewrite step preserves $\Val$, hence the terminal irreducible word has
the same value as $a$.  By Corollary~\ref{cor:fold-as-section}, there is a
unique admissible word with this value, namely $\Fold_\infty(a)$.
\end{proof}

\begin{remark}
Proposition~\ref{prop:rewrite-termination} does not assert confluence of the
abstract rewrite system generated by \textup{(R1)}--\textup{(R4)}.  What is
used here is weaker and sufficient for the paper: every rewrite sequence
terminates, every terminal irreducible is admissible, every rewrite preserves
value, and each value class contains a unique admissible representative.
\end{remark}

For overlays of two admissible inputs the initial alphabet is only
$\{0,1,2\}$, but it is convenient to formulate the rewrite system on the whole
monoid $\D$.  This avoids case distinctions and keeps the normalization target
equal to the canonical section from Corollary~\ref{cor:fold-as-section}.

\subsection{Stable Addition}

For $c,d\in \Xinf$ we define the digitwise superposition
\[
(c+d)_k:=c_k+d_k \in \{0,1,2\}.
\]
Since $c$ and $d$ have finite support, $c+d$ can be viewed as an element of
$\D$.

\begin{definition}[Stable addition]\label{def:stable-add}
For $c,d\in\Xinf$, define
\[
c\oplus d:=Z(\Val(c)+\Val(d)).
\]
\end{definition}

The previous subsection shows that $\oplus$ can be computed without leaving
the digit system.

\begin{theorem}[Addition is superposition followed by normalization]\label{thm:zeckendorf-add-fst}
For all $c,d\in \Xinf$ one has
\[
c\oplus d=\Fold_\infty(c+d).
\]
Equivalently, Zeckendorf addition is the normalization of the noncanonical
digitwise overlay.  Thus the addition problem reduces to the study of the
canonical map $\Fold_\infty$ on the input alphabet $\{0,1,2\}$; the remainder
of the paper compares its one-sided sequential behavior with the classical
most-significant-digit-first Berstel adder.
\end{theorem}

\begin{proof}
By definition of $\Fold_\infty$,
\[
\Fold_\infty(c+d)=Z(\Val(c+d)).
\]
Since $\Val$ is additive on finitely supported digit strings,
\[
\Val(c+d)=\Val(c)+\Val(d).
\]
Therefore
\[
\Fold_\infty(c+d)=Z(\Val(c)+\Val(d))=c\oplus d.
\]
\end{proof}

Once the canonical section has been fixed, the arithmetic problem reduces to
realizing $\Fold_\infty$ by a finite streaming device with bounded delay
\cite{Frougny1992NumbersAutomata,Frougny1999OnlineAddition,
AhlbachUsatineFrougnyPippenger2013}.

\begin{example}[A short normalization trace]\label{rem:worked-normalization-example}
Using the low-to-high Fibonacci-scale convention, let
\[
c=(1,0,1,0,0),\qquad d=(1,0,0,1,0)\in \Xinf.
\]
Then
\[
c+d=(2,0,1,1,0),
\]
so the overlay contains both a digit $2$ and the forbidden pattern $11$.
Its value is
\[
\Val(c+d)=2F_2+F_4+F_5=2+3+5=10=F_3+F_6.
\]
Therefore
\[
\Fold_\infty(c+d)=Z(10)=(0,1,0,0,1)=c\oplus d.
\]
This is the smallest type of example in which normalization must resolve the
two standard obstructions simultaneously.
\end{example}

\section{Unbounded Anticipation and Non-Subsequentiality}\label{sec:streaming}

\subsection{One-Sided Online Delay}

\begin{definition}[One-pass online delay]\label{def:online-delay}
Let $\A_{\mathrm{in}}$ and $\A_{\mathrm{out}}$ be finite alphabets.  A map
\[
T:\A_{\mathrm{in}}^{(\NN)}\to \A_{\mathrm{out}}^{(\NN)}
\]
has \emph{one-pass online delay} $\delta\in\NN$ if, for every $j\ge 1$, the
output digit $(T(x))_j$ depends only on the input prefix
$x_1,\dots,x_{j+\delta}$.  Equivalently, after reading the first $n$ input
digits, a sequential device with delay $\delta$ has already determined the
first $\max\{n-\delta,0\}$ output digits.
\end{definition}

\begin{remark}\label{rem:delay-convention}
Definition~\ref{def:online-delay} is the one-sided low-to-high notion adapted
to $\Fold_\infty$ on finitely supported streams.  It is not the
most-significant-digit-first anticipation parameter used in the classical
online-addition literature \cite{Frougny1999OnlineAddition}.  This distinction
is essential.
\end{remark}

\begin{theorem}[Coordinatewise unbounded anticipation]\label{thm:coordinatewise-unbounded-anticipation}
For every output coordinate $j\ge 1$ and every $\delta\ge 0$, there exist
finitely supported inputs $x,y\in\{0,1,2\}^{(\NN)}$ such that
\[
x_i=y_i \qquad (1\le i\le j+\delta),
\]
but
\[
(\Fold_\infty(x))_j\neq (\Fold_\infty(y))_j.
\]
\end{theorem}

\begin{proof}
Fix $j\ge 1$ and $\delta\ge 0$.

If $j=1$, choose $m\ge 1$ with $2m\ge \delta+1$ and define
\[
x:=\sum_{t=0}^{m-1} e_{2t+1},
\qquad
y:=x+2e_{2m+1}.
\]
The two inputs agree on the first $\delta+1$ positions.  Since $x$ is already
admissible, $\Fold_\infty(x)=x$, so $(\Fold_\infty(x))_1=1$.  Moreover
\[
\Val(x)=\sum_{t=1}^{m}F_{2t}=F_{2m+1}-1,
\]
hence
\[
\Val(y)=F_{2m+1}-1+2F_{2m+2}=F_{2m+4}-1=\sum_{t=1}^{m+1}F_{2t+1}.
\]
Therefore
\[
\Fold_\infty(y)=\sum_{t=1}^{m+1} e_{2t},
\]
and so $(\Fold_\infty(y))_1=0$.

Assume now that $j\ge 2$.  Choose $m\ge 1$ with $2m-1\ge \delta$ and set
\[
x:=\sum_{t=0}^{m-1} e_{j+2t},
\qquad
y:=x+2e_{j+2m}.
\]
The two inputs agree on the first $j+2m-1$ positions, hence on the first
$j+\delta$ positions.  Again $x$ is admissible, so $\Fold_\infty(x)=x$ and
$(\Fold_\infty(x))_j=1$.

For the second input, one computes
\[
\Val(y)
=
\sum_{t=0}^{m-1}F_{j+2t+1}+2F_{j+2m+1}
=
F_{j-1}+\sum_{t=0}^{m}F_{j+2t+2}.
\]
The last identity may be checked directly by induction on $m$ from the
Fibonacci recurrence.
Hence the admissible word
\[
z:=Z(F_{j-1})+\sum_{t=0}^{m} e_{j+1+2t}
\]
has the same value as $y$.  By uniqueness of admissible representatives,
$\Fold_\infty(y)=z$.  Since all nonzero digits of the second summand lie
strictly to the right of position $j$, and $Z(F_{j-1})$ is supported in
positions $<j$, one has $z_j=0$.  Thus $(\Fold_\infty(y))_j=0$.
\end{proof}

\begin{corollary}[No finite one-pass delay]\label{thm:no-bounded-delay}
The canonical normalization map $\Fold_\infty$ admits no finite one-pass
online delay.
\end{corollary}

\begin{proof}
If $\Fold_\infty$ had one-pass online delay $\delta$, then the $j$th output
digit would be determined by the first $j+\delta$ input digits for every
$j\ge 1$.  This contradicts
Theorem~\ref{thm:coordinatewise-unbounded-anticipation}.
\end{proof}

\begin{corollary}[The obstruction already occurs on genuine addition inputs]\label{cor:addition-image-unbounded-anticipation}
For every output coordinate $j\ge 1$ and every $\delta\ge 0$, there exist
admissible words $c,d,c',d'\in \Xinf$ such that the overlays $c+d$ and
$c'+d'$ agree on the first $j+\delta$ coordinates, but
\[
\bigl(c\oplus d\bigr)_j \neq \bigl(c'\oplus d'\bigr)_j.
\]
\end{corollary}

\begin{proof}
In the proof of Theorem~\ref{thm:coordinatewise-unbounded-anticipation}, the
word $x$ is always admissible, so it can be written as $x=x+0$ with
$x,0\in\Xinf$.

If $j=1$, the second witness is
\[
y=\sum_{t=0}^{m-1} e_{2t+1}+2e_{2m+1}
=
\left(\sum_{t=0}^{m} e_{2t+1}\right)+e_{2m+1},
\]
and both summands are admissible.

If $j\ge 2$, the second witness is
\[
y=\sum_{t=0}^{m-1} e_{j+2t}+2e_{j+2m}
=
\left(\sum_{t=0}^{m} e_{j+2t}\right)+e_{j+2m},
\]
and again both summands are admissible.  Hence the witness family already lies
in the image of genuine digitwise superposition of admissible inputs.  The
conclusion then follows from
Theorem~\ref{thm:zeckendorf-add-fst}.
\end{proof}

\subsection{A Finite-Word Consequence}

To compare the one-sided canonical map with standard transducer models, we
encode finitely supported inputs as low-to-high finite words.  For
$u=u_1\cdots u_n\in\{0,1,2\}^\ast$, write
\[
\langle u\rangle:=\sum_{i=1}^n u_i e_i\in \D.
\]
If $a\in\{0,1\}^{(\NN)}$ is finitely supported, let $\operatorname{trim}(a)$
denote the finite word obtained by deleting trailing zeros.  We then define
\[
\mathcal N(u):=\operatorname{trim}(\Fold_\infty(\langle u\rangle))
\qquad (u\in\{0,1,2\}^\ast).
\]

\begin{lemma}[Bounded terminal discrepancy for subsequential transducers]\label{lem:subseq-terminal-bound}
Let $T:\Sigma^\ast\to\Gamma^\ast$ be realized by a deterministic
subsequential transducer with terminal output map $\omega$.  Set
\[
L:=\max_q |\omega(q)|.
\]
If $u$ is a prefix of $v$, then $T(u)$ and $T(v)$ have a common prefix of
length at least $|T(u)|-L$.
\end{lemma}

\begin{proof}
Write $v=uw$.  Let $\alpha$ be the output produced while reading $u$, and let
$q$ be the state reached after $u$.  Then
\[
T(u)=\alpha\,\omega(q).
\]
If the run is continued along $w$, the output has the form
\[
T(v)=\alpha\,\beta\,\omega(q')
\]
for some word $\beta$ and some terminal state $q'$.  Thus $T(u)$ and $T(v)$
share the prefix $\alpha$, whose length is
\[
|\alpha|=|T(u)|-|\omega(q)|\ge |T(u)|-L.
\]
\end{proof}

\begin{theorem}[The one-sided canonical normalizer is not subsequential]\label{thm:canonical-not-subsequential}
The finite-word normalization map
\[
\mathcal N:\{0,1,2\}^\ast\to\{0,1\}^\ast
\]
is not realizable by any deterministic subsequential transducer reading the
input from low Fibonacci scales to high Fibonacci scales.
\end{theorem}

\begin{proof}
Suppose, for contradiction, that a deterministic subsequential transducer
realizes $\mathcal N$, and let $L$ be the constant from
Lemma~\ref{lem:subseq-terminal-bound}.  For each $m\ge 1$, consider the prefix
pair
\[
u_m:=(10)^{m-1}1,
\qquad
v_m:=u_m 02.
\]
Then
\[
\langle u_m\rangle=\sum_{t=0}^{m-1} e_{2t+1},
\qquad
\langle v_m\rangle=\sum_{t=0}^{m-1} e_{2t+1}+2e_{2m+1},
\]
so these are exactly the $j=1$ witnesses from the proof of
Theorem~\ref{thm:coordinatewise-unbounded-anticipation}.  Hence
\[
\mathcal N(u_m)=u_m,
\qquad
\mathcal N(v_m)=0(10)^m1.
\]
In particular, $\mathcal N(u_m)$ and $\mathcal N(v_m)$ differ in their first
symbol, so their common prefix has length $0$.  On the other hand, $u_m$ is a
prefix of $v_m$, and
\[
|\mathcal N(u_m)|=|u_m|=2m-1.
\]
Lemma~\ref{lem:subseq-terminal-bound} therefore yields
\[
0\ge |\mathcal N(u_m)|-L = 2m-1-L
\]
for every $m$, which is impossible for $m>L$.  This contradiction shows that
$\mathcal N$ is not subsequential.
\end{proof}

\begin{theorem}[Exact prefix-destruction index]\label{thm:exact-prefix-destruction}
For $n\ge 1$, define
\[
\Delta(n):=
\max\Bigl\{
n-\lcp(\mathcal N(u),\mathcal N(v)):
u\text{ is a proper prefix of }v,\ |\mathcal N(u)|=n
\Bigr\}.
\]
Then
\[
\Delta(n)=n
\qquad (n\ge 1).
\]
\end{theorem}

\begin{proof}
Since $\lcp(\mathcal N(u),\mathcal N(v))\ge 0$, one always has
$\Delta(n)\le n$.

If $n=2m-1$ is odd, use the prefix pair
\[
u_m=(10)^{m-1}1,
\qquad
v_m=u_m02
\]
from the proof of Theorem~\ref{thm:canonical-not-subsequential}.  There
\[
\mathcal N(u_m)=u_m,
\qquad
\mathcal N(v_m)=0(10)^m1,
\]
so $|\mathcal N(u_m)|=2m-1=n$ and
$\lcp(\mathcal N(u_m),\mathcal N(v_m))=0$.  Hence $\Delta(2m-1)=2m-1$.

If $n=2m$ is even, set
\[
a_m:=(10)^{m-1}2,
\qquad
b_m:=a_m1.
\]
Then $a_m$ is a proper prefix of $b_m$.  Moreover
\[
\Val(a_m)=\sum_{t=1}^{m-1}F_{2t}+2F_{2m}
=(F_{2m-1}-1)+2F_{2m}=F_{2m+2}-1
=\sum_{t=1}^{m}F_{2t+1},
\]
so
\[
\mathcal N(a_m)=(01)^m.
\]
Likewise,
\[
\Val(b_m)=\Val(a_m)+F_{2m+1}=F_{2m+3}-1
=\sum_{t=1}^{m+1}F_{2t},
\]
and therefore
\[
\mathcal N(b_m)=1(01)^m.
\]
Thus $|\mathcal N(a_m)|=2m=n$ and
$\lcp(\mathcal N(a_m),\mathcal N(b_m))=0$, so $\Delta(2m)=2m$.

Both parity classes attain the trivial upper bound, hence $\Delta(n)=n$ for
all $n\ge 1$.
\end{proof}

\begin{remark}\label{rem:frougny-separation}
The point of Theorem~\ref{thm:canonical-not-subsequential} is not to claim a
first non-sequentiality phenomenon in Fibonacci numeration.  Rather, it gives
a sharpened statement for the specific canonical low-to-high normalization map:
the obstruction is coordinatewise, it is witnessed by an explicit family of
prefix pairs, and it already appears on genuine addition overlays.
\par
Theorem~\ref{thm:canonical-not-subsequential} shows that the one-sided map
$\Fold_\infty$ is not the same object as the classical online adders of
Frougny and Berstel.  In Frougny's most-significant-digit-first model,
addition in base $\varphi$ is computable with delay $3$, and the associated
finite-word recoding is subsequential, but the output is not the greedy
expansion \cite[Prop.~13]{Frougny1999OnlineAddition}.  The next two sections
revisit that classical transducer on its own terms.
\end{remark}

\section{The Classical Berstel Adder}\label{sec:kernel}

\subsection{A Deterministic Subsequential Transducer}

We now switch to the classical finite-word, most-significant-digit-first
setting of \cite{Frougny1999OnlineAddition,LabbeLepsova2023FibonacciTwosComplement}.
The relevant device is the $10$-state Berstel adder.  Because the final output
depends on the last state, the natural model is a deterministic
\emph{subsequential} transducer rather than a letter-to-letter Mealy machine.

Its state set is
\[
Q=
\{000,001,002,010,100,101,0\bar{1}2,1\bar{1}2,01\bar{1},11\bar{1}\}.
\]
The transitions are
\[
\begin{aligned}
000&\xrightarrow{d/0}00d,\qquad 100\xrightarrow{d/1}00d\quad(d=0,1,2),\\
001&\xrightarrow{0/0}010,\ \xrightarrow{1/0}100,\ \xrightarrow{2/0}101,\\
002&\xrightarrow{0/0}11\bar{1},\ \xrightarrow{1/1}000,\ \xrightarrow{2/1}001,\\
010&\xrightarrow{0/0}100,\ \xrightarrow{1/0}101,\ \xrightarrow{2/1}0\bar{1}2,\\
101&\xrightarrow{0/1}010,\ \xrightarrow{1/1}100,\ \xrightarrow{2/1}101,\\
0\bar{1}2&\xrightarrow{0/0}01\bar{1},\ \xrightarrow{1/0}010,\ \xrightarrow{2/0}100,\\
1\bar{1}2&\xrightarrow{0/1}01\bar{1},\ \xrightarrow{1/1}010,\ \xrightarrow{2/1}100,\\
01\bar{1}&\xrightarrow{0/0}001,\ \xrightarrow{1/0}002,\ \xrightarrow{2/0}1\bar{1}2,\\
11\bar{1}&\xrightarrow{0/1}001,\ \xrightarrow{1/1}002,\ \xrightarrow{2/1}1\bar{1}2.
\end{aligned}
\]
The terminal output map is
\[
\omega(000)=000,\ \omega(001)=001,\ \omega(002)=010,\ \omega(010)=010,
\]
\[
\omega(100)=100,\ \omega(101)=101,\ \omega(0\bar{1}2)=000,\ \omega(1\bar{1}2)=100,
\]
\[
\omega(01\bar{1})=001,\qquad \omega(11\bar{1})=101.
\]

Let
\[
\mathcal K=(Q,\delta,\eta,\omega)
\]
denote the deterministic subsequential transducer defined above, with initial
state $000$.  For a state $q\in Q$, let
\[
F_q:\{0,1,2\}^\ast\to\{0,1\}^\ast
\]
be the function obtained by starting in $q$, following the transition labels,
and appending the terminal word $\omega$ of the last state.
For a finite word $w=d_1\cdots d_n$ read most significant digit first, we
write
\[
\Val_{\mathrm{MSD}}(w):=\sum_{i=1}^{n} d_i F_{n-i+2}.
\]
When discussing residual functions and shortest separating words, we keep the
complete output word, including leading zeros, because those zeros belong to
the subsequential output.  They are immaterial only at the level of the
represented Fibonacci value.

All ten displayed states are reachable from the initial state.  One may use
\[
\varepsilon,\ 1,\ 2,\ 10,\ 11,\ 12,\ 20,\ 102,\ 202,\ 1020
\]
to reach respectively
\[
000,\ 001,\ 002,\ 010,\ 100,\ 101,\ 11\bar{1},\ 0\bar{1}2,\ 1\bar{1}2,\ 01\bar{1}.
\]

\begin{proposition}[Classical correctness]\label{prop:classical-berstel-correctness}
The transducer $\mathcal K$ is the classical Berstel adder in the
most-significant-digit-first Fibonacci setting.  It computes a value-preserving
binary recoding of words over $\{0,1,2\}$: for every input word $w$, the
complete output $\mathcal K(w)$ is binary and satisfies
\[
\Val_{\mathrm{MSD}}(\mathcal K(w))=\Val_{\mathrm{MSD}}(w).
\]
In general this output is not the greedy Zeckendorf expansion.
\end{proposition}

\begin{proof}
This is the classical construction; see
\cite[Prop.~13]{Frougny1999OnlineAddition} and
\cite[Thm.~3.1]{LabbeLepsova2023FibonacciTwosComplement}.  Those references
show that the machine realizes a value-preserving recoding from ternary
Fibonacci input words to binary Fibonacci output words in the
most-significant-digit-first convention.  The non-greedy statement is already
part of Frougny's proposition.
\end{proof}

\begin{example}[A short non-greedy output witness]\label{ex:berstel-nongreedy}
Starting from the initial state $000$, the input word $121$ produces the
complete output word $001100$.  After suppressing the two leading zeros, the
represented binary word is $1100$, which still contains the forbidden factor
$11$.  Hence the classical Berstel output is not the greedy Zeckendorf
expansion of the same value.
\end{example}

\begin{remark}
Theorem~\ref{thm:no-bounded-delay} and
Proposition~\ref{prop:classical-berstel-correctness} concern different
sequential functions.  The former concerns the canonical one-sided map
$\Fold_\infty$, whereas the latter concerns the classical Berstel recoding in
the most-significant-digit-first model.  Example~\ref{ex:berstel-nongreedy}
exhibits the difference
explicitly: the Berstel machine returns a non-greedy value-preserving recoding,
while the canonical map of Section~4 is defined to return the greedy
representative.
\end{remark}

\section{Minimality of the Berstel Kernel}\label{sec:minimality}

\subsection{Residual Functions}

Define an equivalence relation on $Q$ by
\[
q\sim q' \iff F_q=F_{q'}.
\]
For deterministic subsequential transducers, this is the usual sequential
Nerode equivalence.  Any deterministic subsequential transducer implementing
the same function as $\mathcal K$ must have at least $|Q/\!\sim|$ reachable
states.

To avoid ambiguity, we distinguish the displayed realization $\mathcal K$ from
the subsequential function that it induces.  The reachability statement above
shows that $\mathcal K$ really has ten reachable states as a concrete machine.
The minimality problem addressed below is a statement about the induced
function: namely, whether any deterministic subsequential realization of that
same function can use fewer than ten reachable states.

\begin{lemma}[Separation by terminal output]\label{lem:terminal-separation}
If $\omega(q)\neq \omega(q')$, then $q\not\sim q'$.  Consequently the only
pairs not already separated by the empty input are
\[
(000,0\bar{1}2),\quad
(001,01\bar{1}),\quad
(002,010),\quad
(100,1\bar{1}2),\quad
(101,11\bar{1}).
\]
\end{lemma}

\begin{proof}
If $\omega(q)\neq \omega(q')$, then
\[
F_q(\varepsilon)=\omega(q)\neq \omega(q')=F_{q'}(\varepsilon),
\]
so $q\not\sim q'$.  Reading the terminal output list above, one sees that the
only equal-terminal pairs are exactly the five pairs displayed.
\end{proof}

\begin{lemma}[The remaining pairs are separated by the input $0$]\label{lem:zero-separation}
The five pairs listed in Lemma~\ref{lem:terminal-separation} are all
inequivalent.  More precisely,
\[
F_{000}(0)=0000 \neq 0001 = F_{0\bar{1}2}(0),
\]
\[
F_{001}(0)=0010 \neq 0001 = F_{01\bar{1}}(0),
\]
\[
F_{002}(0)=0101 \neq 0100 = F_{010}(0),
\]
\[
F_{100}(0)=1000 \neq 1001 = F_{1\bar{1}2}(0),
\]
\[
F_{101}(0)=1010 \neq 1001 = F_{11\bar{1}}(0).
\]
\end{lemma}

\begin{proof}
Each identity is obtained by following one transition on the input $0$ and
then appending the terminal output of the arrival state.
\end{proof}

\subsection{Minimality of the Classical Kernel}

\begin{theorem}[Minimality of the Berstel kernel]\label{thm:berstel-kernel-minimality}
All states in $Q$ are reachable and pairwise inequivalent under subsequential
Nerode equivalence.  Consequently the function induced by $\mathcal K$ has
deterministic subsequential state complexity $10$, and $\mathcal K$ is a
minimal realization of that function.
\end{theorem}

\begin{proof}
For deterministic subsequential transducers, two states are equivalent if and
only if every input word produces the same complete output word from both
states, including the terminal output.  Thus the theorem reduces to separating
the residual functions $F_q$.

By Lemma~\ref{lem:terminal-separation}, every pair with distinct terminal
output is already separated by the empty input.  Lemma~\ref{lem:zero-separation}
handles the only five remaining pairs.  Hence all ten states are pairwise
inequivalent, and in fact every inequivalent pair is separated by an input
word of length at most $1$.  Equivalently, the partition-refinement
computation reaches the discrete partition
\[
\boxed{|Q/\!\sim|=10,\qquad
\max_{q\neq q'}|w_{q,q'}|\le 1,}
\]
where $w_{q,q'}$ denotes a shortest separating word for the pair $(q,q')$.
Hence every state defines a distinct residual sequential function, so no
smaller deterministic subsequential transducer can realize the same map.  As
all ten states are reachable, the displayed transducer $\mathcal K$ itself is
therefore a minimal realization.
\end{proof}

\begin{corollary}[Exact state complexity]\label{cor:berstel-state-complexity}
The classical Berstel adder has deterministic subsequential state complexity
$10$.
\end{corollary}

\begin{proof}
This is exactly Theorem~\ref{thm:berstel-kernel-minimality}.
\end{proof}

\begin{corollary}\label{cor:berstel-kernel-unique}
The minimal deterministic subsequential realization of the function induced by
$\mathcal K$ is unique up to labeled isomorphism.
\end{corollary}

\begin{proof}
This is the standard uniqueness statement for minimal deterministic sequential
machines; see \cite{BerstelReutenauer1988RationalSeries}.
\end{proof}

\section{A Local Parallel Lower Bound}\label{sec:complexity}

The failure of bounded one-pass delay does not, by itself, rule out efficient
synchronous local normalization.  The next theorem shows that locality also
fails in that model: the number of rounds must grow linearly with the input
scale.

\begin{theorem}[Local parallel lower bound for Zeckendorf normalization]\label{thm:zeckendorf-parallel-propagation-lower-bound}
Let $\A=\{0,1,2\}$ and define, for finitely supported one-sided configurations
$x=(x_1,x_2,\ldots)\in\A^{(\NN)}$,
\[
\Val(x):=\sum_{k\ge 1} x_kF_{k+1}.
\]
Let $Z(\Val(x))\in\{0,1\}^{(\NN)}$ denote the unique admissible Zeckendorf
representative with the same value.

Suppose $\Phi:\A^{(\NN)}\to\{0,1\}^{(\NN)}$ is a synchronous local map of
radius $r$, meaning that after extending every input by zeros outside $\NN$
there exists a local rule $f:\A^{2r+1}\to\{0,1\}$ such that
\[
(\Phi(x))_j=f(x_{j-r},\dots,x_{j+r})
\qquad (j\ge 1),
\]
and let $\Phi^t$ denote its $t$-fold iterate.  If for some $t\ge 1$ one has
\[
\Phi^t(x)=Z(\Val(x))
\]
for every finitely supported input $x$, then for every even integer $m\ge 2$,
\[
t\ge \left\lceil \frac{m-1}{r}\right\rceil.
\]
In particular, no fixed-radius synchronous local normalization scheme can run
in a number of rounds independent of the input scale.
\end{theorem}

\begin{proof}
After $t$ iterations, the output coordinate at position $j$ depends only on
the input in the interval $[j-rt,j+rt]\cap\NN$.
Fix an even integer $m\ge 2$ and define
\[
x^{(m)}_k=
\begin{cases}
1,&k\in\{1,3,5,\dots,m-1\},\\
0,&\text{otherwise},
\end{cases}
\qquad
y^{(m)}:=x^{(m)}+e_1.
\]
The configuration $x^{(m)}$ is already admissible, hence
$Z(\Val(x^{(m)}))=x^{(m)}$ and its $m$th coordinate is $0$.

On the other hand,
\[
\Val(x^{(m)})
=
\sum_{j=1}^{m/2}F_{2j}
=
F_{m+1}-1,
\]
so
\[
\Val(y^{(m)})=\Val(x^{(m)})+F_2=F_{m+1}.
\]
Therefore $Z(\Val(y^{(m)}))$ is the admissible word with a single $1$ at
position $m$, and hence
\[
Z(\Val(x^{(m)}))_m=0,
\qquad
Z(\Val(y^{(m)}))_m=1.
\]

If $rt<m-1$, then the coordinate $(\Phi^t(\cdot))_m$ cannot depend on the input
at position $1$, so it must take the same value on $x^{(m)}$ and $y^{(m)}$.
This contradicts the required normalization outputs above.  Hence
$rt\ge m-1$, which is equivalent to the stated bound.
\end{proof}

\begin{remark}
Theorem~\ref{thm:zeckendorf-parallel-propagation-lower-bound} concerns the
nonredundant input alphabet $\{0,1,2\}$ and the canonical Zeckendorf output
map $x\mapsto Z(\Val(x))$.  It does not contradict constant-time parallel
addition results for redundant digit sets, where carries can be absorbed by
extra local freedom.  We also do not claim a matching upper bound here:
the theorem only shows that every fixed-radius synchronous realization of the
canonical map requires a number of rounds growing linearly with the input
scale.
\end{remark}

\section{Conclusion}

We have separated three layers of the Zeckendorf addition problem.  First, the
canonical normalization map $\Fold_\infty$ is the section of the Fibonacci
congruence selected by admissible representatives, and it admits a direct
rewrite-theoretic description.  Second, in the one-sided low-to-high model,
this canonical map has coordinatewise unbounded anticipation and therefore no
bounded online delay; more strongly, its low-to-high finite-word incarnation
is not subsequential.  Third, the classical $10$-state Berstel adder belongs
to a different, most-significant-digit-first subsequential setting, and in
that setting its state complexity is exactly $10$.  The local lower bound
shows that canonical normalization also remains nonlocal in the synchronous
setting.  Natural next questions concern analogous distinctions for other
Pisot numeration systems and for redundant or signed digit alphabets.
Machine-checked verification of the rewrite system and of the minimal
subsequential quotient of the classical kernel would also be worthwhile.

\bibliographystyle{abbrv}
\bibliography{references_godel_zeckendorf,references_fibonacci,references}

\end{document}
